\documentclass[12pt]{article}
\usepackage[utf8]{inputenc}
\usepackage{amsmath, amssymb, amsthm}
\usepackage{geometry}
\usepackage{graphicx}
\usepackage{listings}
\usepackage{color}
\usepackage{float}
\usepackage{enumitem}

\geometry{margin=1in}

% Code snippet styling
\definecolor{codegray}{rgb}{0.5,0.5,0.5}
\definecolor{codepurple}{rgb}{0.58,0,0.82}
\definecolor{backcolour}{rgb}{0.95,0.95,0.92}

\lstdefinestyle{mystyle}{
    backgroundcolor=\color{backcolour},   
    commentstyle=\color{codegray},
    keywordstyle=\color{magenta},
    numberstyle=\tiny\color{codegray},
    stringstyle=\color{codepurple},
    basicstyle=\ttfamily\footnotesize,
    breakatwhitespace=false,         
    breaklines=true,                 
    captionpos=b,                    
    keepspaces=true,                 
    numbers=left,                    
    numbersep=5pt,                  
    showspaces=false,                
    showstringspaces=false,
    showtabs=false,                  
    tabsize=2
}
\lstset{style=mystyle}

\title{Homework 1 Solutions}
\author{Aniket Ghosh \\ STAT 6530}


\begin{document}

\maketitle

\section*{Problem 1}
\textbf{Consider the function $g(y)$ defined as:}
\[ g(y) = \begin{cases} y^2, & \text{for } 0 < y < 1, \\ 0, & \text{elsewhere.} \end{cases} \]

\subsection*{(i) Is $g(y)$ a pdf?}
To check if the function is a pdf, we integrate it over its support to check if it equals 1.
\[ \int_{0}^{1} y^2 \, dy = \left[ \frac{y^3}{3} \right]_0^1 = \frac{1}{3} - 0 = \frac{1}{3} \]
Since $\frac{1}{3} \neq 1$, $g(y)$ is not a pdf.

\subsection*{(ii) Find the constant $c$ so that $cg(y)$ is a pdf. Denote this pdf by $f(y)$.}
We need $\int_{0}^{1} c \cdot y^2 \, dy = 1$. Since we know that $\int_{0}^{1} y^2 \, dy = \frac{1}{3}$, we can take $c$ to be $3$.
Thus, the pdf is:
\[ f(y) = 3y^2, \quad 0 < y < 1 \]

\subsection*{(iii) Show that the cdf is given by $F(y) = y^3$ for $0 < y < 1$.}
The cdf is defined as $F(y) = \int_{-\infty}^{y} f(t) \, dt$.
For $0 < y < 1$:
\[ F(y) = \int_{0}^{y} 3t^2 \, dt = \left[ t^3 \right]_0^y = y^3 \]

\subsection*{(iv) Finding the inverse function $F^{-1}(p)$ and the three percentiles.}
$p = F(y) = y^3$. Solving for $y$:
\[ y = p^{1/3} \]
Thus, $F^{-1}(p) = p^{1/3}$ for $0 < p < 1$.

\noindent The percentiles are calculated as follows:
\begin{itemize}
    \item \textbf{ $p=0.25 = (0.25)^{1/3} \approx 0.62996$}
    \item \textbf{ $p=0.50 = (0.50)^{1/3} \approx 0.79370$}
    \item \textbf{ $p=0.75 = (0.75)^{1/3} \approx 0.90856$}
\end{itemize}

\subsection*{(v) Simulation using Inverse CDF method in R}
\begin{lstlisting}[language=R]
set.seed(123)
# Generate uniform random variables
random_var <- runif(10000)
# Inverse CDF transformation
inverse_f <- random_var^(1/3)

# Calculate sample quartiles
sample_quartiles <- quantile(inverse_f, probs = c(0.25, 0.50, 0.75))
print(sample_quartiles)
\end{lstlisting}

\noindent \textbf Using this R code, we obtain the values: 0.6323802, 0.7908156, and 0.9058849, respectively. We see that the sample quartiles from the simulation are very close to the theoretical values calculated.

\subsection*{(vi) Relative Frequency Histogram}
\begin{lstlisting}[language=R]
hist(inverse_f, prob=TRUE, main="Histogram of Simulated Y vs PDF of f(y)", 
     xlab="y", col="lightblue", breaks=50)
curve(3*x^2, from=0, to=1, col="red", lwd=2, add=TRUE)
legend("topleft", legend="Theoretical PDF 3y^2", col="red", lwd=2)
\end{lstlisting}

\begin{figure}[H]
    \centering
    % Make sure the filename inside {} matches exactly what you uploaded
    \includegraphics[width=0.8\textwidth]{histogram_1.png}
    \caption{Relative frequency histogram of simulated data overlaid with PDF of $f(y)$.}
\end{figure}

\noindent \textbf The relative frequency histogram of the simulated data closely matches the pdf curve $f(y) = 3y^2$, showing an exponential increase in density as $y$ approaches 1.

\newpage

\section*{Problem 2}
\textbf{Suppose $Y$ is a continuous random variable with the pdf $f_Y(y)$ and the cdf $F_Y(y)$. For any number $c$, $R(c) = \int_{c}^{\infty} f_Y(t)dt = 1 - F_Y(c)$. Consider the random variable $R(Y) = 1 - F_Y(Y)$.}

\subsection*{(a) Find the cdf of $R(Y)$.}
Let $F_{R(Y)}(r)$ denote the cumulative distribution function of the random variable $R(Y)$. By definition:
\[ F_{R(Y)}(r) = P(R(Y) \le r) \]
Substituting the expression for $R(Y)$ given in the problem:
\[ F_{R(Y)}(r) = P(1 - F_Y(Y) \le r) \]
Rearranging the inequality:
\[ F_{R(Y)}(r) = P(F_Y(Y) \ge 1 - r) \]
We know that for any continuous random variable $Y$, the transformation $F_Y(Y)$ follows a standard Uniform distribution on the interval $(0, 1)$. 
For a Uniform(0,1) variable, the probability of being greater than or equal to some value $k$ is $1 - k$ (provided $0 < k < 1$).
Therefore, substituting $1-r$ for $k$:
\[ F_{R(Y)}(r) = 1 - (1 - r) = r \]
Thus, the cdf of $R(Y)$ is:
\[ F_{R(Y)}(r) = \begin{cases} 0, & r \le 0 \\ r, & 0 < r < 1 \\ 1, & r \ge 1 \end{cases} \]

\noindent This is a uniform distribution itself.

\subsection*{(b) Find the pdf of $R(Y)$.}
Let $f_{R(Y)}(r)$ denote the probability density function of $R(Y)$. The pdf is the derivative of the cdf with respect to $r$: 
\[ f_{R(Y)}(r) = \frac{d}{dr} F_{R(Y)}(r) = \frac{d}{dr}(r) = 1 \]
Thus, the pdf is:
\[ f_{R(Y)}(r) = \begin{cases} 1, & 0 < r < 1 \\ 0, & \text{elsewhere} \end{cases} \]
This also indicates that $R(Y)$ follows a Uniform(0,1) distribution.

\subsection*{(c) Find the mean and the variance of $R(Y)$.}
Using the pdf, we can calculate the expected value (mean) and variance directly.

\noindent \textbf{Mean:}
\[ E[R(Y)] = \int_{-\infty}^{\infty} r \cdot f_{R(Y)}(r) \, dr = \int_0^1 r \cdot 1 \, dr = \left[ \frac{r^2}{2} \right]_0^1 = \frac{1}{2} \]

\noindent \textbf{Variance:}
First, we calculate the second moment $E[R(Y)^2]$:
\[ E[R(Y)^2] = \int_0^1 r^2 \cdot f_{R(Y)}(r) \, dr = \int_0^1 r^2 \, dr = \left[ \frac{r^3}{3} \right]_0^1 = \frac{1}{3} \]
Now, using the variance formula $Var(R(Y)) = E[R(Y)^2] - (E[R(Y)])^2$:
\[ Var(R(Y)) = \frac{1}{3} - \left(\frac{1}{2}\right)^2 = \frac{1}{3} - \frac{1}{4} = \frac{1}{12} \]

\newpage

\section*{Problem 3}
\textbf{Suppose $Y$ is a standard normal random variable with cdf $\Phi(y) = P[Y \le y] = \int_{-\infty}^{y} \frac{1}{\sqrt{2\pi}} e^{-t^2/2} dt$.}

\subsection*{(a) Show that for any $c$, $\Phi(c) + \Phi(-c) = 1$.}
By definition, $\Phi(c) = P(Y \le c)$ and $\Phi(-c) = P(Y \le -c)$.
Because the standard normal probability density function $\phi(t) = \frac{1}{\sqrt{2\pi}}e^{-t^2/2}$ is symmetric about 0 (because it has a distribution $N~(0,1$), the area under the curve to the left of $-c$ is equal to the area under the curve to the right of $c$.
Mathematically:
\[ P(Y \le -c) = P(Y \ge c) \]
We also know that the total probability is 1:
\[ P(Y \le c) + P(Y \ge c) = 1 \]
Substituting $\Phi(-c)$ for $P(Y \ge c)$ (using the symmetry property derived above):
\[ \Phi(c) + \Phi(-c) = 1 \]

\subsection*{(b) Let $X = \Phi(-Y)$. Find the cdf of $X$ and the pdf of $X$.}
First, we simplify the expression for $X$ using the identity derived in part (a):
\[X = \Phi(-Y) = 1 - \Phi(Y) \]


\noindent To find the cdf $F_X(x) = P(X \le x)$, we substitute our simplified $X$:
\[ F_X(x) = P(1 - \Phi(Y) \le x) \]
Rearranging the inequality:
\[ F_X(x) = P(\Phi(Y) \ge 1 - x) \]
It is a known property that for any continuous random variable $Y$ with cdf $\Phi$, the random variable $\Phi(Y)$ follows a standard Uniform distribution on $(0,1)$.
For a standard uniform variable, the probability of being greater than or equal to a value $k$ is simply $1-k$. Here, $k = 1-x$.
Therefore:
\[ F_X(x) = 1 - (1 - x) = x \]
Thus, the cdf is:
\[ F_X(x) = \begin{cases} 0, & x \le 0 \\ x, & 0 < x < 1 \\ 1, & x \ge 1 \end{cases} \]

\noindent To find the pdf $f_X(x)$, we take the derivative of the cdf with respect to $x$:
\[ f_X(x) = \frac{d}{dx} F_X(x) = \frac{d}{dx}(x) = 1 \]
Thus, the pdf is:
\[ f_X(x) = \begin{cases} 1, & 0 < x < 1 \\ 0, & \text{elsewhere} \end{cases} \]

\subsection*{(c) Find the mean and the variance of $X$.}
We see that the distribution is a $U~(0,1)$ distribution, and we have already calculated the mean and variance in the last problem.

\noindent\textbf{Mean: $\frac{1}{2}$}

\noindent \textbf{Variance: $\frac{1}{12}$}

\newpage

\section*{Problem 4}
\textbf{Consider two random variables $X$ and $Y$.}

\subsection*{(a) Show that $Var(X+Y) = Var(X) + Var(Y) + 2Cov(X,Y)$.}
By definition of variance:
\[ Var(X+Y) = E[((X+Y) - E[X+Y])^2] \]
Let $\mu_x = E[X]$ and $\mu_y = E[Y]$. Then $E[X+Y] = \mu_x + \mu_y$.
\begin{align*}
Var(X+Y) &= E[((X - \mu_x) + (Y - \mu_y))^2] \\
&= E[(X - \mu_x)^2 + (Y - \mu_y)^2 + 2(X-\mu_x)(Y-\mu_y)] \\
&= E[(X - \mu_x)^2] + E[(Y - \mu_y)^2] + 2E[(X-\mu_x)(Y-\mu_y)] \\
&= Var(X) + Var(Y) + 2Cov(X,Y)
\end{align*}

\subsection*{(b) Show that $Var(X-Y) = Var(X) + Var(Y) - 2Cov(X,Y)$.}
Using the result from (a) for random variables $X$ and $-Y$:
\[ Var(X + (-Y)) = Var(X) + Var(-Y) + 2Cov(X, -Y) \]
We know that $Var(-Y) = (-1)^2 Var(Y) = Var(Y)$ and $Cov(X, -Y) = -Cov(X,Y)$.
Substituting these in:
\[ Var(X-Y) = Var(X) + Var(Y) - 2Cov(X,Y) \]

\subsection*{(c) Show how (a) and (b) simplify when X and Y are uncorrelated.}
If $X$ and $Y$ are uncorrelated, then $Cov(X,Y) = 0$.
The equations simplify to:
\[ Var(X+Y) = Var(X) + Var(Y) \]
\[ Var(X-Y) = Var(X) + Var(Y) \]

\newpage

\section*{Problem 5}
\textbf{Suppose $Y$ has an exponential distribution with pdf $f(y;\lambda) = \lambda e^{-\lambda y}$ and cdf $F(y;\lambda) = 1 - e^{-\lambda y}$.}

\subsection*{(a) Find the median of Y.}
At the median, $F(y) = 0.5$:
\[ 1 - e^{-\lambda y} = 0.5 \implies e^{-\lambda y} = 0.5 \]
Take $\ln$ of both sides:
\[ -\lambda y = \ln(0.5) = -\ln(2) \implies y = \frac{\ln(2)}{\lambda} \approx \frac{0.693}{\lambda} \]

\subsection*{(b) Find the lower quartile and upper quartile.}
\textbf{Lower Quartile ($Q_1$):} Set $F(q_1) = 0.25$.
\[ 1 - e^{-\lambda q_1} = 0.25 \implies e^{-\lambda q_1} = 0.75 \]
\[ q_1 = \frac{-\ln(0.75)}{\lambda} = \frac{\ln(4/3)}{\lambda} \approx \frac{0.288}{\lambda} \]

\noindent \textbf{Upper Quartile ($Q_3$):} Set $F(q_3) = 0.75$.
\[ 1 - e^{-\lambda q_3} = 0.75 \implies e^{-\lambda q_3} = 0.25 \]
\[ q_3 = \frac{-\ln(0.25)}{\lambda} = \frac{\ln(4)}{\lambda} \approx \frac{1.386}{\lambda} \]

\subsection*{(c) Comparing median and upper quartile with the mean.}
The mean of an exponential distribution is $E[Y] = \frac{1}{\lambda}$.
\begin{itemize}
    \item \textbf{Median vs Mean:} $\frac{0.693}{\lambda} < \frac{1}{\lambda}$. The median is smaller than the mean.
    \item \textbf{Upper Quartile vs Mean:} $\frac{1.386}{\lambda} > \frac{1}{\lambda}$. The upper quartile is larger than the mean.
\end{itemize}

\newpage

\section*{Problem 6}
\textbf{In an exit poll, $n=1648$, $\hat{p} = 0.515$ for Candidate A.}

\subsection*{(a) Find the standard error assuming $p=0.50$.}
The standard error of the sample proportion is $\sqrt{\frac{p(1-p)}{n}}$.
Using $p=0.50$:
\[ SE = \sqrt{\frac{0.50(1-0.50)}{1648}} = \sqrt{\frac{0.25}{1648}} \approx 0.0123 \]

\subsection*{(b) Simulation Results}

\textbf{(i) Simulation Code (R):}
\begin{lstlisting}[language=R]
set.seed(123)
n_trials <- 100000
n_voters <- 1648
# Simulate election under assumption p = 0.5
simulations <- rbinom(n_trials, n_voters, 0.5) / n_voters

# Calculate probability of observing 0.515 or higher
prob_extreme <- mean(simulations >= 0.515)
print(paste("Probability of observing >= 51.5%:", prob_extreme))
\end{lstlisting}

\noindent We get the following results: $P(X \geq 51.5\%) = 0.11336$. We see that there is a 11.336$\%$ chance of this result. 

\noindent \textbf{(ii) Using the Standard Error}
We know that for a normal distribution, approximately 95\% of the data lies within 2 standard deviations (standard errors) of the mean.
We calculate this "non-surprising" range around the assumed population mean of $0.50$:
\[ \text{Range} = 0.50 \pm 2(SE) \]
\[ \text{Range} = 0.50 \pm 2(0.0123) = 0.50 \pm 0.0246 \]
\[ \text{Interval} = [0.4754, 0.5246] \]
The observed exit poll proportion of 0.515 falls comfortably inside this interval. Therefore, obtaining a result of 51.5\% when the true population is split 50/50 is not surprising, because the result is within the range of expected variability. Hence, it would be difficult to justify predicting the results of this election based on this poll.

\newpage

\section*{Problem 7}

\section*{Problem 7}

Consider a sequence $X_1, X_2, \dots$ of independent $Ber(p)$ random variables. Let $T_n = \frac{1}{n}\sum_{i=1}^n X_i$. We define the logit function $g(x) = \log\left(\frac{x}{1-x}\right)$.
We want to find the asymptotic distribution of:
\[ \sqrt{n}\left[g(T_n) - g(p)\right] \]

Since each $X_i \sim Ber(p)$, we have $E[X_i] = p$ and $Var(X_i) = p(1-p)$. According to the Central Limit Theorem, the standardized quantity converges to a standard normal distribution:
\[ \sqrt{n}(T_n - p) \xrightarrow{d} N(0, p(1-p)) \]

We apply the Delta Method, which states that for a sequence of random variables $T_n$ satisfying $\sqrt{n}(T_n - p) \xrightarrow{d} N(0, \sigma^2)$, the transformed variable converges as follows:
\[ \sqrt{n}(g(T_n) - g(p)) \xrightarrow{d} N\left(0, [g'(p)]^2 \sigma^2\right) \]

We must find the derivative $g'(p)$. Rewriting the function using logarithm rules simplifies the differentiation:
\[ g(x) = \log\left(\frac{x}{1-x}\right) = \log(x) - \log(1-x) \]
Differentiating with respect to $x$:
\[ g'(x) = \frac{d}{dx}(\log(x)) - \frac{d}{dx}(\log(1-x)) = \frac{1}{x} - \left(\frac{1}{1-x} \cdot (-1)\right) = \frac{1}{x(1-x)}  \]

Evaluated at $p$, we have $g'(p) = \frac{1}{p(1-p)}$.

Finally, we compute the asymptotic variance by substituting $g'(p)$ and $\sigma^2 = p(1-p)$ into the Delta Method formula:
\[ Var_{asymp} = [g'(p)]^2 \cdot \sigma^2 = \left[ \frac{1}{p(1-p)} \right]^2 \cdot p(1-p) \]
\[ Var_{asymp} = \frac{p(1-p)}{p^2(1-p)^2} = \frac{1}{p(1-p)} \]

Thus, the final convergence in distribution is:
\[ \sqrt{n}\left[\log\left(\frac{T_n}{1-T_n}\right) - \log\left(\frac{p}{1-p}\right)\right] \xrightarrow{d} N\left(0, \frac{1}{p(1-p)}\right) \]

\newpage

\section*{Problem 8}
\textbf{Using Chebyshev's inequality to prove the law of large numbers.}

Chebyshev's inequality states that if $X$ is a random variable with finite variance $\sigma^2$ and mean $\mu$, then for any $k > 0$:
\[ P(|X - \mu| \ge k\sigma) \le \frac{1}{k^2} \]

Let $X_1, X_2, \dots, X_n$ be i.i.d. random variables with mean $\mu$ and finite variance $\sigma^2$.
Let $\overline{X}_n = \frac{1}{n}\sum_{i=1}^n X_i$.
The expected value and variance of the sample mean are:
\[ E[\overline{X}_n] = \mu \]
\[ Var(\overline{X}_n) = \frac{\sigma^2}{n}, \quad \text{so } SD(\overline{X}_n) = \frac{\sigma}{\sqrt{n}} \]

We want to show that for any $\epsilon > 0$:
\[ \lim_{n \to \infty} P(|\overline{X}_n - \mu| \ge \epsilon) = 0 \]

Applying Chebyshev's inequality to the random variable $Y = \overline{X}_n$. We set the distance $\epsilon$ equal to $k$ standard deviations:
\[ \epsilon = k \cdot SD(\overline{X}_n) = k \frac{\sigma}{\sqrt{n}} \]
Solving for $k$:
\[ k = \frac{\epsilon\sqrt{n}}{\sigma} \]
Substituting this into Chebyshev's bound ($\frac{1}{k^2}$):
\[ P(|\overline{X}_n - \mu| \ge \epsilon) \le \frac{1}{(\frac{\epsilon\sqrt{n}}{\sigma})^2} = \frac{\sigma^2}{n\epsilon^2} \]
Taking the limit as $n \to \infty$:
\[ \lim_{n \to \infty} P(|\overline{X}_n - \mu| \ge \epsilon) \le \lim_{n \to \infty} \frac{\sigma^2}{n\epsilon^2} = 0 \]
Since probabilities cannot be negative, the probability converges to 0. This proves the Law of Large Numbers.

\end{document}